\chapter{PELAKSANAAN PENELITIAN}
\label{pelaksanaan-penelitian}


\section{Alat dan Bahan Penelitian}


Pada bagian ini:
\begin{enumerate}[label=\alph*)]
\item Uraikan secara rinci spesifikasi dan jangkauan kemampuan alat yang digunakan. Alat
bisa berupa perangkat keras (hardware) maupun perangkat lunak (software).
\item Jika penelitian melibatkan penggunaan bahan-bahan (kimiawi, fisik, dll.), uraikan
spesifikasi bahan yang digunakan.
\item Jika penelitian bersifat empirik, gambarkan rancangan sistem alat untuk penelitian.
\end{enumerate}



\section{Tata Laksana Penelitian}
 Uraikan rangkaian logis penyelesaian masalah menurut tahap-tahap
analisis yang dipaparkan dalam bagian Dasar Teori, yaitu berupa langkah-langkah
kerja dan/atau algoritma penelitian.


\section{Rencana Analisis Hasil}
 Kemukakan bagaimana, menurut rencana, hasil-hasil yang akan
diperoleh dari penelitian akan diolah. Cara bagaimana pengolahan ini akan dilakukan
sudah tentu disesuaikan/dikaitkan dengan tujuan penelitian. Secara umum, pengolahan bisa
dilakukan melalui proses:
\begin{enumerate}[label=\alph*)]
\item Perangkuman hasil penelitian dalam format tabel, gambar, statistik (rata-rata,
koefisien korelasi, dlsb.), atau dalam bentuk besaran khusus tertentu sesuai dengan
parameter atau variabel yang dilibatkan dalam penelitian.
\item Pengujian perbedaan statistik (rata-rata, korelasi, dlsb) variabel penelitian.
\item Pengujian keterkaitan (korelasi) statistik variabel penelitian.
\item Pengolahan lain yang relevan dengan tujuan penelitian.
\end{enumerate}


